\chapter{Interoperabilidad con SIG}\label{anexo-g.-interoperabilidad-gis}

Este anexo detalla la interoperabilidad con herramientas SIG resumida en la sección~\ref{exportacion-a-geotiff}: el formato de exportación, el estilo de calado y el flujo de trabajo para abrir los resultados en QGIS.

\section{Cloud-Optimized GeoTIFF}\label{g-cog}

El formato de intercambio es el \textit{Cloud-Optimized GeoTIFF} (COG), un GeoTIFF que organiza la imagen en tiles internos y añade vistas reducidas (\textit{overviews}) precalculadas. Esta organización permite leer solo la región y el nivel de detalle necesarios en lugar del fichero entero, la misma filosofía de lectura parcial que motiva el uso de TileDB (sección~\ref{rendimiento-consultas}). Al ser un GeoTIFF corriente, cualquier herramienta que lea GeoTIFF lo abre. El exportador lo genera con la opción \texttt{-{}-cog}, marcando el fichero como \textit{tiled} con bloques de 512$\times$512 píxeles, construyendo \textit{overviews} en los factores 2, 4, 8 y 16 por submuestreo promediado y manteniendo la compresión DEFLATE sin pérdida.

\section{Estilo de calado}\label{g-qml}

Junto al COG se genera un fichero de estilo QML (el formato de estilos de QGIS) que QGIS aplica automáticamente al cargar la capa. Usa una rampa de color por umbrales de calado coherente con la aplicación web, con las fronteras derivadas de Russo et al.\ \cite{russo2013} recogidas en la Tabla~\ref{tab:qml-stops}. De este modo, un mapa abierto en QGIS y otro generado por la aplicación se interpretan con el mismo criterio cromático.

\begin{table}[ht]
\centering
\small
\begin{tabularx}{\textwidth}{>{\raggedright\arraybackslash}p{3.0cm}>{\raggedright\arraybackslash}p{3.0cm}L}
\toprule
\textbf{Umbral de calado} & \textbf{Color} & \textbf{Categoría} \\
\midrule
$\geq$ 0,01 m & azul muy claro & Húmedo \\
0,25 m & azul claro & Peligroso para niños \\
0,50 m & azul medio & Peligroso para adultos \\
1,00 m & azul oscuro & Crítico \\
$\geq$ 2,00 m & azul muy oscuro & Extremo \\
\bottomrule
\end{tabularx}
\caption{Paradas de color del estilo QML de calado, alineadas con los umbrales de peligrosidad de Russo et al.\ \cite{russo2013}.}
\label{tab:qml-stops}
\end{table}

\section{Flujo de trabajo}\label{g-flujo}

Todo el proceso se encapsula en un único comando que, dado un escenario y un instante temporal, genera el paquete completo (COG, estilo QML y proyecto de QGIS):

\begin{quote}
\texttt{python gis\_demo.py -{}-dataset datos1 -{}-step 10\_00 -{}-out paquete\_gis}
\end{quote}

El proyecto \texttt{.qgz} referencia el COG mediante ruta relativa, de modo que el paquete puede entregarse como una carpeta autocontenida sin que se rompan los enlaces: quien lo recibe solo tiene que abrir el proyecto en QGIS y la capa aparece ya estilizada, con la misma codificación de peligrosidad que la aplicación web.
